% CFG and SSA value flow / 控制流图与 SSA 值流
\begin{figure}[htbp]
  \centering
  \begin{tikzpicture}[x=.94cm,y=1cm]
    \node[vis/artifact,minimum width=2.7cm] (entry) at (0,1.55)
      {\texttt{entry}\quad $i_0=0,\ acc_0=0$};
    \node[vis/semantic node,text width=4.35cm] (loop) at (0,0)
      {\texttt{loop}\quad $i=\phi(0,i')$\\
       $acc=\phi(0,acc')$\quad $i<k$?\\
       $I(i):\ acc=\sum_{j<i}c_j\ \wedge\ 0\le i\le k$};
    \node[vis/artifact,text width=3.55cm] (body) at (-1.15,-2.35)
      {\texttt{body}\quad load--multiply--clip\\$i'=i+1,\quad acc'=acc+term$};
    \node[vis/artifact,minimum width=2.25cm] (exit) at (1.9,-2.35)
      {\texttt{exit}\\return $acc$};

    \draw[vis/edge] (entry) -- node[vis/label,right]{初值} (loop);
    \draw[vis/edge] (loop.south west) -- node[vis/label,left]{true} (body.north);
    \draw[vis/edge] (loop.south east) -- node[pos=.58,vis/label,right]{false} (exit.north);
    \draw[vis/edge] (body.west) to[out=180,in=180,looseness=1.45]
      node[vis/label,left]{回边} (loop.west);
    \draw[vis/semantic edge] (body.east) to[out=12,in=-25,looseness=1.25]
      (loop.east);

    \begin{scope}[on background layer]
      \node[vis/group,fit=(entry)(loop)(body)(exit),inner sep=6pt] (cfg) {};
    \end{scope}
    \node[vis/label,font=\bfseries\footnotesize,above=2mm of cfg] {控制流与值流};

    \node[vis/decision,text width=3.65cm,anchor=north west] (corr) at (4.0,1.85)
      {$i$：缩放地址 $\Rightarrow$ 指针递增\\
       $term$：比较/选择 $\Rightarrow$ 寄存器值\\
       $acc$：跨回边值 $\Rightarrow$ \texttt{addw}};
    \node[vis/note,text width=3.65cm,below=4mm of corr] (why)
      {黑色实线决定下一基本块；蓝色粗线只追踪回边携带的值。$\phi$ 的选择取决于实际到达的前驱边。};
    \draw[vis/deferred edge] ([yshift=-1mm]cfg.north east)
      to[out=20,in=165] node[pos=.52,vis/label,above]{降低时对应}
      (corr.north west);
  \end{tikzpicture}
  \caption{循环同时是一张控制流图和一组跨回边的数据依赖：$\phi$ 节点表达“来自哪条前驱边的值”，而后端可以把索引归纳改写为指针递增，只要这些依赖保持不变。}
  \label{fig:ssa-cfg}
\end{figure}
